Notre préoccupation, lors de l'élaboration de ce projet de cartographie du traitement documentaire a donc été, non pas de proposer un outil pour l'exploration des collections étant donné qu'ils existent déjà et que de nouveaux outils sont développés, mais d'un outil proposant aux utilisateurs de nouvelles perspectives d'accessibilité aux bases de données qui permettent de contourner les contraintes d'accessibilité dues à l'évolution des politiques documentaires. Cet outil a été pensé en complémentarité aux outils de consultation des données, comme un outil améliorant l'accessibilité les collections afin de permettre aux utilisateurs de préparer leur recherches. 

Dans ce but, nous allons déjà présenter les outils dont nous disposons pour améliorer la découvrabilité des collections par le traitement documentaire. Puis nous conclurons notre propos en évoquant comment nous pouvons adapter le processus de recherche des utilisateurs de la base de données à ces enjeux d'accessibilité notamment en appuyant l'importance que nous avons porté à l'expérience utilisateur dans le cadre d'une recherche avec notre outil.  